\section{Related Work}\label{sec:related}

The DDP~\cite{ddp} model wrapper, which is based on the model replication design, was an initial distributed training feature introduced in PyTorch~\cite{pytorch:nips:2019}. Although it can handle large datasets, it cannot accommodate the ever-increasing model sizes that are now prevalent in the field.

ZeRO~\cite{zero, zero:offload} and cross-replica sharding~\cite{xu2020automatic} inspired the FSDP design, but FSDP is intrinsically different. Prior work employs model partitioning or per-parameter sharding to distribute parameter tensors, and rely on \code{Broadcast} and \code{Gather} collective communication primitives to synchronize values. Although this design can achieve the same functionality, it could lead to uneven workload distribution across GPU devices, which hampers the efficiency of synchronized distributed training. Additionally, since this approach modifies the internals of the machine learning framework, such as tensor storage and memory management, it might no longer work when the internal implementation is updated or new features are introduced. Therefore, a native solution that is co-designed with the core components of the framework would provide a more robust and consistent user experience.


MiCS~\cite{zhang2022mics} and FSDP differ in gradient communication strategies. MiCS uses a global \code{AllReduce} followed by sharding within each partition group, whereas FSDP employs \code{AllGather} and \code{ReduceScatter}. As a result, each rank in MiCS must hold the entire model gradients, leading to higher memory usage than FSDP's approach of sharding a single layer. While both MiCS and FSDP use a hybrid communication strategy to improve efficiency at scale, FSDP's approach schedules \code{AllGather} within a flexibly-sized sharded group, potentially resulting in lower runtime latency than the two-hop \code{AllGather} utilized by MiCS. This reduced latency is crucial as the \code{AllGather} operation is critical to execution, and limiting the world size and participants of \code{AllGather} to accelerators within a group with good locality can result in lower latency and higher throughput.

Pipeline parallelism~\cite{gpipe, harlap2018pipedream} involves partitioning model parameters and their activations across multiple devices through the division of models into pipeline stages. However, pipeline parallelism requires model changes and meticulous tuning for microbatch sizes, number of stages and partitions, as well as intricate scheduling procedures to optimize performance by squeezing out bubbles. 

Additionally, specific attention is given to high profile architectures such as transformers. For example, sequence parallelism~\cite{sequence} reduces activation memory in conjunction with tensor parallelism; Pipetransformer~\cite{pipetransformer} designed a dynamic 2D parallelism that allows changing the dimensions of pipeline and data parallelism on the fly, depending on learning signals. These methods are highly effective but can be difficult to generalize as they either rely on the specific implementation or the model's layered structure.

Many existing solutions combine data parallelism with other parallelisms to achieve speedup. For example, Megatron~\cite{narayanan2021efficient} demonstrated highly efficient deep transformer training on large clusters using 3D (data, tensor and pipeline) parallelism. Further, compiler-based techniques such as Alpa~\cite{zheng2022alpa}, GSPMD~\cite{gspmd}, and FlexFlow~\cite{flexflow} leverage profiling, performance modeling, user annotations and search to find the best configuration across the parallelism space of data, tensor and pipeline for a given cluster. In all cases, FSDP provides the benefit of being a drop-in replacement for data parallelism that reduces data redundancy along the data parallel axis.



Orthogonal memory-saving techniques include gradient compression~\cite{compression}, mixed-precision training~\cite{precision}, tensor rematerialization~\cite{DTR} and CPU-offloading~\cite{fang2022parallel}, but they could have implications on model accuracy and incur overhead in (un)compression, quantization, recomputation, and host-to-device copies, respectively.